\hypertarget{project-kennel-design-principles}{%
\chapter{Project Kennel: Design
Principles}\label{project-kennel-design-principles}}

These are the axioms the design and architecture derive from, sorted in
four small categories: Premise, System, Architecture and Boundaries.
They state posture, not mechanism.

Every structural choice in the framework answers to one or more of them;
these are cited throughout to justify the mechanisms that realise them.

\hypertarget{premise}{%
\section{Premise}\label{premise}}

\hypertarget{good-boy}{%
\subsection{\texorpdfstring{good-boy}{good-boy}}\label{good-boy}}

``Good boy'' is praise for intention, and the principle is that
intention is not what the boundary judges. The name describes a
disposition --- loyal, well-meaning, eager to please --- and a
disposition is the one thing confinement refuses to consult. What a
workload intends and what it does are different things, and the boundary
answers only to the second; to confine by intention is to confuse the
dog's character with the state of the couch. The framework never asks
whether a workload means well, because meaning well is not what bounds
the damage.

A boundary that holds only for the workloads you correctly judged
untrustworthy fails the moment the judgment is off --- and the judgment
is always eventually off. The insider was fine until they weren't; the
allowlisted dependency shipped clean releases until the one that wasn't;
the agent was loyal until a docstring in the repo redirected it. Kennel
takes the judgment out of the loop. The living room is protected from
any dog, so you are never breached by having misjudged one.

What a workload can reach and do is the whole of the damage; what it
meant changes nothing. An AI agent reading every file ``in case
something is useful'' and a credential-harvesting postinstall script
perform the same reads and leave the same hole, and the boundary treats
them identically because intent does not move the blast radius. The dog
loves you and wrecks the couch anyway. You do not crate it as a verdict
on its character; you keep the couch safe from whatever the dog turns
out to be. This is why the project is confinement and not detection ---
you are never required to be right about the workload, and so being
wrong about it is not a breach.

\hypertarget{split-the-uid}{%
\subsection{\texorpdfstring{split-the-uid}{split-the-uid}}\label{split-the-uid}}

Unix fuses who you are with what you may do. The uid is both your
identity and your authority --- your shell, your editor, your compiler
are you, and reach everything you reach by the number alone. That fusion
was sound while the only code running as you was code acting for you,
and it is what broke when the code became an agent, a postinstall
script, a freshly cloned build: still your uid, no longer you.

Kennel does not answer this by giving the workload a different identity.
It keeps the user's uid --- deliberately, because a separate identity
would demand the delegated id-ranges and privilege the framework is
built to do without --- and splits the authority off the identity
instead. The uid stays; the ambient reach it would have carried does
not. In its place the framework defines the context the uid runs in, the
view and the channels and the signed policy, so the workload holds not
the user's inherited authority but exactly the authority the operator
declared. Identity and authority, fused by Unix, come apart into two
contexts under one uid --- the default, where authority is ambient and
total, and the kennel, where it is explicit and bounded. Without the
split the monitor has nothing to mediate, since an un-split uid carries
its ambient authority straight through any policy laid over it.

\hypertarget{do-less}{%
\subsection{\texorpdfstring{do-less}{do-less}}\label{do-less}}

Confronted with a workload whose intent it will not judge and whose
authority it has had to strip, the framework's founding instinct is to
do as little as possible. Every line of mediation is a line that can be
wrong; every component is a thing that must be trusted and read; every
active mechanism is a surface an adversary can work against. So the
question asked of each part of the design is not ``how do we do this
well'' but ``can we avoid having to do it at all'' --- and the strongest
answer to how a thing is secured is usually that the thing is not there
to secure. This is the keystone, and most of the architecture is its
consequences. \texttt{construction-by-absence} is do-less carried to its
limit against a resource --- remove it, and there is nothing left to get
right. \texttt{control-not-data-plane} is do-less on the byte path, the
monitor absent from the flow it authorises. \texttt{rule-of-1} is
do-less inside each component, which carries one hazard rather than
three. The reference monitor is verifiable, Anderson's third property,
because a monitor that does less is small enough to read in full. The
framework grows only by necessity argued in the open, and the answer
that removes a thing is worth more than the answer that improves it.

\hypertarget{system}{%
\section{System}\label{system}}

\hypertarget{reference-monitor}{%
\subsection{\texorpdfstring{reference-monitor}{reference-monitor}}\label{reference-monitor}}

A confinement system is a reference monitor or it is scenery. Anderson
named the three properties such a monitor must hold: it mediates every
access to a protected resource and cannot be bypassed; the code it
constrains cannot alter it, nor can any policy it loads lower the floor
it guarantees, however that policy was signed; and it is small enough to
be analysed for correctness. These are one test to pass, not three
features to collect --- a monitor that meets two and misses the third
does not confine. The host level grew such a monitor over decades, in
mandatory access control, the LSM framework, and systemd's syscall and
address-family filters; the user level, where the workloads that now
matter actually run, never did. Kennel is that monitor, assembled from
existing kernel primitives and run in user space, where every access
passes through it, the workload it confines can neither bypass nor alter
it, and it is kept small enough to read in full. Its three properties
organise the register: complete mediation is reached two ways, by
\texttt{construction-by-absence} and \texttt{interpose-as-transaction};
tamperproofing is why the enforcer is never placed inside the thing it
confines, and why it re-asserts its own invariants at every spawn
whoever signed the policy, since a monitor that could be signed out of
its guarantees would cut the floor from beneath its own feet; and
verifiability is \texttt{do-less} made a requirement, with
\texttt{rule-of-1} and \texttt{control-not-data-plane} its sharpest
applications. \texttt{refuse-to-start} is what keeps a degraded monitor
from passing as one.

\hypertarget{mediate-use-not-reach}{%
\subsection{\texorpdfstring{mediate-use-not-reach}{mediate-use-not-reach}}\label{mediate-use-not-reach}}

Kernel access control answers one question --- may this channel be
opened --- and then stands down. It never sees the method call on the
bus it allowed, the destination behind the socket it permitted, the
secret read through the file it granted. The platform conceded this long
ago and patched around it one protocol at a time: dbus-daemon filters
individual method calls, polkit gates privileged actions, ssh forced
commands and git-shell constrain a single connection's use. Each exists
because the channel was the wrong granularity. A sandbox that confines
by reachability inherits the same ceiling --- once a socket or a path is
present in the workload's world, what the workload does through it is
ungoverned. Kennel mediates the use rather than the reach. The content
of the interaction is what policy decides, and a granted channel is the
start of mediation, not the end of it --- a monitor that stopped at the
channel it allowed would be \texttt{reference-monitor} in name only.

The principle is the design's default and not an absolute. One rung is
deliberately exempt: where a protocol's whole authority is exactly what
the operator meant to grant, the channel may be handed over raw,
mediation ending at the connection because there is no smaller unit of
use worth governing. That exemption is per-service and
operator-declared, loud and threat-tagged, never the silent default ---
the law is mediate the use, and the raw grant is the named place the
operator chooses to stop at the reach and own what follows.

\hypertarget{construction-by-absence}{%
\subsection{\texorpdfstring{construction-by-absence}{construction-by-absence}}\label{construction-by-absence}}

Complete mediation (\texttt{reference-monitor}) is reached two ways. The
first makes the resource absent, so there is no access to mediate. If a
resource is not present to the workload, the count of accesses to it is
zero and complete mediation holds with nothing to enforce. So Kennel
does not hand the workload the real world and subtract --- it constructs
the workload's view of each resource class from nothing, admitting only
what policy grants. The filesystem holds the granted paths and no
others; an unreachable network is empty, not blocked; a withheld socket
has no name in the workload's world to connect to. What is absent is not
denied on access --- it is not present, not deniable, not enumerable. A
workload cannot probe for what is not there, cannot map the user's
environment by reading the edges of its permissions, cannot tell a
withheld resource from one that never existed. Three things follow that
a denylist cannot offer: inspection is one-directional, since ``what can
the workload see'' is the constructed view itself, with no
everything-not-denied left to reason about; policy edits are positive, a
grant adding to the view rather than a deny chasing what was forgotten;
and the view is the boundary, so an authoring error is a visible
inclusion, not an invisible omission. Absence is a stronger mediation
than denial, because there is nothing left to attack --- it is
\texttt{do-less} carried to its limit, the resource removed rather than
guarded.

\hypertarget{interpose-as-transaction}{%
\subsection{\texorpdfstring{interpose-as-transaction}{interpose-as-transaction}}\label{interpose-as-transaction}}

Some resources cannot be made absent --- the workload genuinely needs to
reach a service, a destination, a socket --- and for these the first
limb, \texttt{construction-by-absence}, has nothing to remove. The
second way to complete mediation is interposition: the workload does not
hold the resource, it holds a request for it, and every use becomes a
transaction the monitor authorises one call at a time. A granted network
destination is not an open socket the workload owns but a connect the
monitor performs on its behalf and hands back; a granted service is not
a path in the view but a brokered call decided per invocation. Because
the decision and the action are a single authorised transaction, there
is no window between checking a request and honouring it for the
workload to widen --- freedom from time-of-check/time-of-use races falls
out of the structure rather than being guarded for. Which limb applies
to a resource is a deliberate per-class choice: remove it where you can,
interpose where you cannot, and complete mediation holds either way.

\hypertarget{architecture}{%
\section{Architecture}\label{architecture}}

\hypertarget{rule-of-1}{%
\subsection{\texorpdfstring{rule-of-1}{rule-of-1}}\label{rule-of-1}}

A confinement system is defined by where it concentrates risk. The
common heuristic is the Rule of 2: a component should not combine more
than two of three hazards --- it parses untrusted input, it is written
in a memory-unsafe language, or it runs with privilege. Kennel holds its
trusted base to a stricter line, a Rule of 1, where no component carries
more than one of the three at once. Where \texttt{do-less} removes
components from the base, this keeps the survivors simple: the part that
decodes a workload's bytes runs unprivileged and memory-safe; the part
that holds privilege parses nothing a workload influenced, and parses it
before any sandbox exists that could shape its input. The baseline
language is safe Rust, and where \texttt{unsafe} is required to drive a
kernel primitive --- the syscalls and ioctls the sandbox is built from
--- it is quarantined into a component that neither holds privilege nor
parses untrusted input. A bug in any one component is then bounded by
that component's single hazard, never compounded into the
untrusted-input-and-privilege-and-unsafe-memory combination from which
escapes are built.

\hypertarget{control-not-data-plane}{%
\subsection{\texorpdfstring{control-not-data-plane}{control-not-data-plane}}\label{control-not-data-plane}}

The monitor decides; it does not carry. When a workload is granted a
channel, the daemon authorises and establishes the connection, then
steps out of the byte path entirely --- the established descriptor is
handed to the facade, and the data flows without passing through the
daemon again. This is a performance property, since steady-state I/O is
never bottlenecked on a central mediator, but it is first a security
one: a component that never sees a byte of workload traffic cannot be
exploited by that traffic. Keeping the daemon on the control plane keeps
adversarial input out of the most-trusted, most-reachable component in
the system. This is \texttt{do-less} on the byte path --- the monitor
made absent from the flow it authorises rather than secured within it,
the same removal \texttt{rule-of-1} makes inside each component. The
monitor is in the loop for every decision and out of the loop for every
byte.

\hypertarget{authentication-never-attestation}{%
\subsection{\texorpdfstring{authentication-never-attestation}{authentication-never-attestation}}\label{authentication-never-attestation}}

A kennel exists to run code the operator has chosen not to trust, and
that choice has a hard edge: the workload may prove an identity, but it
may never vouch as the operator. The distinction is between
authentication and attestation. Authentication is host-verifiable and
can be mediated --- the framework can let a workload prove the user's
identity to a specific destination while binding which destination, so
the credential is used but never aimed by the workload. Attestation is
the operator putting their name on data --- a signed commit, a signed
release, a gpg signature --- and it cannot be delegated to an entity the
operator does not trust, because there is no host-side safety property
to bind it to: the signature is over content the workload authored, and
a signature an untrusted workload can produce is the operator's identity
stamped on whatever that workload chose. So the framework brokers
authentication and refuses attestation outright. The human signs on
review, after seeing what they are signing; the workload never holds the
pen.

\hypertarget{boundaries}{%
\section{Boundaries}\label{boundaries}}

\hypertarget{request-dont-author}{%
\subsection{\texorpdfstring{request-dont-author}{request-dont-author}}\label{request-dont-author}}

The workload asks; it does not write. Every capability a workload can
exercise traces to a policy the operator signed, and the workload's role
is to request what that policy already allows, never to author the
policy itself or invent a capability the operator did not grant. This
holds even where the workload supplies input to the decision --- an
agent that may spawn another kennel selects from within a signed
template's declared bounds, filling fields the template opened, but the
floor it spawns against is the operator's signature, not the agent's
request. The separation is the whole point: an actor that could author
its own grants would be confined by nothing, since the boundary would
move wherever the workload pushed it. Consistent with \texttt{good-boy},
authority is never extended to a workload on the strength of trusting
it, because trust is not a thing the framework weighs;
request-don't-author keeps the boundary where the operator put it, and
makes the workload's influence a choice among permitted options rather
than a power to permit.

\hypertarget{footgun-warn-dont-forbid}{%
\subsection{\texorpdfstring{footgun-warn-dont-forbid}{footgun-warn-dont-forbid}}\label{footgun-warn-dont-forbid}}

The operator is the trust root, and a framework that forbids what the
operator has legitimately decided to do is not protecting them --- it is
substituting its judgement for theirs and inviting them to route around
it. So where a capability degrades security but is conceivably the right
call --- a media kennel that needs a mount service, a query to a
connectivity daemon --- the default-deny is overridable in a delta, with
a loud warning and a threat tag that makes the deviation visible instead
of silent. The cost is shown; the choice stays the operator's. The
principle is bounded by its opposite, and the boundary is narrow: a
small, named set of capabilities are refused outright rather than
warned, because permitting them would be theatre --- claiming
confinement while handing over a signing oracle or a spawn escape that
empties the claim. That is the same refusal \texttt{refuse-to-start}
makes when the substrate cannot back a policy: warn for the footgun the
operator might reasonably fire, but refuse the one whose only effect is
to make the confinement a lie.

\hypertarget{tell-me-why}{%
\subsection{\texorpdfstring{tell-me-why}{tell-me-why}}\label{tell-me-why}}

Every grant carries the reason it exists. A delta that adds a path,
opens a destination, or removes a default is not accepted without a
stated justification, and that justification travels with the grant
wherever it goes: into the diff a reviewer reads, into the settled
policy's provenance, into the audit record of the decision the grant
authorised. A grant without a reason is one no one can review on the day
it is written or understand six months later, and the framework refuses
to accept one. The reason does not gate the grant; an operator with
standing to make a change makes it. What the reason does is make the
change accountable, so that a confinement stays legible to the people
who must trust it rather than becoming a wall of permissions whose
origins have been forgotten. Where \texttt{footgun-warn-dont-forbid}
makes a security-degrading grant loud, this makes every grant
answerable, the loud warning being the sharp end of a discipline that
runs through all of them. A policy no one can read is a policy no one
can review, and a confinement no one reviews drifts from the floor it
started above.

\hypertarget{refuse-to-start}{%
\subsection{\texorpdfstring{refuse-to-start}{refuse-to-start}}\label{refuse-to-start}}

A policy is a claim about what a workload cannot do, and a claim the
substrate cannot enforce is worse than no claim, because the operator
believes it. If a kennel's policy needs a kernel mechanism the host does
not provide --- the network isolation, the execution gate, the socket
policy a given grant depends on --- Kennel refuses to start it rather
than run it degraded and silent. A monitor that runs without the
mechanism its policy depends on is precisely the scenery
\texttt{reference-monitor} opens by warning against: present, plausible,
and enforcing nothing. Silent degradation is the specific failure this
forbids --- a kennel that asked to block lateral movement and then ran
without the mechanism that blocks it has not failed safe, it has lied to
the operator while appearing to work. Failing closed turns an
unenforceable policy into a visible error the operator must resolve, and
restores the one property that makes a confinement claim worth anything:
that when it says a thing is prevented, the thing is prevented.
